\section{2. Use case and building context}

\textbf{The building.} BuildSim models the LTU A-house: three floors and 838
rooms, with a REST API for equipment, sensor values and actuator states, and a
3D viewer. It stores the current state only. It has no physics, no history, no
clock, and it comes back blank after a restart. We instrument 50 rooms on the
ground floor: the 3 fika rooms, 20 lecture rooms, 24 offices and 3 small rooms
that the people simulation never uses. They are a seeded random sample, listed
in \code{deploy/devices.csv}, so the set is the same in every run.

\textbf{The people.} occupancysim, a course example, generates each day from a
seed and the date: staff in their offices, fika breaks around 09:45 and 14:30,
lunch at 12:00--13:00, lectures in four fixed slots, and three guards on patrol
at night. It also owns the model clock, which runs 60 or 120 times faster than
real time in our runs. Every time and duration in this document is model time
unless it says ``real''.

\textbf{How our processes use BuildSim.} The devices write to BuildSim, and
only the devices:
\begin{itemize}
\tightlist
\item Each \textbf{sensor device} writes its reading
  (\code{PUT /api/sensors/\{id\}/value}).
\item Each \textbf{actuator device} writes the state it has reached
  (\code{PUT /api/actuators/\{id\}/state}).
\item \textbf{physics} reads the actuators' reached states from BuildSim and
  never writes to it.
\item The \textbf{registry} creates the equipment records, and recreates them
  when BuildSim restarts.
\item \textbf{viz} writes the room layers and alerts the viewer shows.
\end{itemize}
This differs from the report guide, which has the simulator write sensor values
into BuildSim and sensor processes read them back. We follow Lectures 1 and 2
instead: sensor processes produce the readings, and the true room state never
enters BuildSim or anything the controller can read. The course material
disagrees with itself here, and the lectures' version is the one that keeps the
truth out of the controller's reach.

\textbf{Actors.} The \emph{facility manager} watches air quality, energy and
alerts, in the 3D viewer and on our graphs page. The \emph{installer} mounts
devices and records each one's id, kind and room in the registry; a device does
nothing until that is done. \emph{Occupants} affect the system only by being in
a room. None of them is on the control path.

\textbf{Why the decision is not trivial.} Air quality and energy pull in
opposite directions, and a room answers slowly. Its \COtwo{} time constant is
the room volume divided by the airflow: about 12 minutes for a lecture room at
design flow, about an hour for an office, and about two hours for any room
while its damper sits at the occupied minimum. A controller that waits for the
\COtwo{} sensor to show that a room has filled is already behind. One that
knows the room is about to fill can open the damper early and gently. Whether
that is worth anything here is what the evaluation in Section 11 measures, by
comparing four policies over the same simulated week: constant design flow,
reactive control on measured \COtwo{}, our predictive policy, and an oracle
that knows the true future occupancy.

\textbf{Scenario for the demonstration.} On a weekday morning the fika room
1570 is empty and its \COtwo{} is at 417~ppm. At 09:28 the predictive policy
raises its airflow to 112~l/s for the 11 people the profile expects within the
room's lead time of 22 minutes; the break starts at about 09:45. Then a
\COtwo{} sensor is frozen: after five identical readings the controller stops
believing it, commands design flow, and the dashboard raises an alert. Finally
the controller is stopped: every damper's last command expires after 300
seconds and it returns to design flow on its own.

\todo{Confirm that this is the demo you will give, or replace it. The first two
steps have been run (2026-09-30 smoke test; T-36). The third, stopping the
controller, has no recorded run yet (T-10, T-23).}
